\section{Discussion}
\label{sec:discussion}

\subsection{Why $b \gg \rtip$}

The large offset $b = \SI{1979}{\um}$ compared to the pole tip radius
$\rtip = \SI{5}{\um}$ is a direct consequence of the yokeless magnetic
circuit. Without a ferromagnetic yoke to guide the flux from one pole tip
to another, the return path passes entirely through air, which has a much
higher reluctance than iron. The total reluctance
$\Ra = 3.19 \times 10^8$ A/Wb includes:

\begin{itemize}
\item $R_\mathrm{pole}$: reluctance of the soft-iron core (small);
\item $R_\mathrm{gap}$: reluctance of the air gap between tip and sample (small);
\item $R_\mathrm{return\;air}$: reluctance of the air return path (dominant).
\end{itemize}

Because $R_\mathrm{return\;air}$ is large, the flux does not concentrate
sharply at the tip but instead spreads out as if originating from a virtual
charge far behind the surface. This is the physical origin of the large $b$.

In contrast, Menq's designs~\cite{Zhang2010,Menq2011} use a ferromagnetic
yoke that provides a low-reluctance return path, effectively short-circuiting
$R_\mathrm{return\;air}$. This forces the flux to concentrate at the tips,
yielding $b = l = \SI{405}{\um}$---nearly five times smaller.

\subsection{Reluctance Comparison Caveats}

As noted in \cref{sec:derived-params}, the $\Ra$ values from this work and
from Menq are not directly comparable because they represent different
quantities:
\begin{itemize}
\item \textbf{This work}: $\Ra = R_\mathrm{total}$ (single pole, no yoke).
\item \textbf{Menq}: $\Ra = R_\mathrm{air\text{-}gap}$ only (yoke provides
  $R_\mathrm{return} \to 0$).
\end{itemize}
The factor-of-5 difference in $\Ra$ (this work is \emph{smaller}) does not
mean the yokeless design is magnetically more efficient; rather, it reflects
the different definitions. The offset $b$ is the only parameter that provides
a fair, topology-independent comparison.

\subsection{Distance Range Dependence of $b$}

The sensitivity of $b$ to the fitting distance range
(Section~\ref{sec:range-sensitivity}) is physically meaningful. At short
distances, higher-order multipole terms---which decay faster than
$1/(x+b)^2$---contribute to the measured field. Including only short-range
data causes the fit to absorb these contributions into a smaller effective
$b$. The full-range estimate ($b = \SI{1979}{\um}$) represents the true
monopole component; the short-range estimate ($b \approx \SI{1504}{\um}$)
is biased by dipole and quadrupole corrections.

This behavior could be exploited in future work: fitting with a
monopole-plus-dipole model would separate the two contributions and provide
additional information about the pole geometry.

\subsection{Implications for Actuator Design}

The results have clear implications for the design of magnetic force actuators:

\begin{enumerate}
\item \textbf{Yoke is essential for strong forces.} The yokeless configuration
produces forces approximately $240\times$ weaker than a yoke-based design at
comparable distances, making it unsuitable for applications requiring
nanonewton-scale forces.

\item \textbf{$b$ as a design metric.} The offset parameter $b$ provides a
single number that captures the effectiveness of flux concentration,
independent of coil turns, current, or sensor sensitivity. It can be used
to compare actuators across different geometries and topologies.

\item \textbf{Frequency-domain calibration is robust.} The excellent frequency
independence ($\mathrm{CV}_a = 0.06\%$, $\mathrm{CV}_b = 0.06\%$) validates
the use of low-frequency transfer function measurements for magnetic
characterization. This approach avoids the complications of DC offset drift
and provides intrinsic common-mode rejection.
\end{enumerate}
